\section{Discussion and Conclusion}\label{sec:conclusion}
This paper presents CODA, which formulates image-based score following as
cascaded selection over known candidates and introduces a silence-driven
break mode for jump recovery. We also contribute a repeat-aware jump test benchmark with manually annotated repeat structures for the MSMD test set, providing the first standardized evaluation protocol for discontinuity handling in image-based score following. Despite achieving state-of-the-art performance, several limitations remain. CODA is trained and evaluated exclusively on solo piano from the MSMD dataset~\cite{dorfer2018learning}. Larger datasets that pair real recordings with scanned scores would benefit not only this work but the field as a whole. Extending CODA to other solo instruments is a natural first step. Beyond solo, chamber music and orchestral settings pose additional challenges due to timbral overlap, denser layouts, and multiple simultaneous parts. Finally, evaluation under real-world conditions, including scanned pages, commercial recordings, and live microphone input, is needed to further assess the method.

